\SourceLeaf{363}%
Kreises durch \(AB\), ferner eine Linie \(=AC-AB\) durch \(CF\) bezeichnet, durch Schlüsse in der bei den alten Geometern üblichen Form folgende Proportionen ab:
\par
{\fontsize{10}{12}\selectfont
\[
\begin{aligned}
AC:AB&=5AB+2CF:3AB+CF=19AB+7CF:11AB+4CF\\
&=71AB+26CF:41AB+15CF=265AB+97CF:153AB+56CF\\
&=989AB+362CF:571AB+209CF
\end{aligned}
\]
}
Aus der vorletzten folgt dann leicht \(AC:AB>265:153\), so wie aus der letzten, wenn man eine Linie \(BD=2AB-AC=AB-CF\) einführt,
\[
AC:AB=1351AB-362BD:780AB-209BD<1351:780
\]
Dass Hr. M., welcher sich mit der bei den alten Geometern üblichen Einkleidung arithmetischer Schlüsse sehr vertraut gemacht hat, \textsc{Archimed}'s Ideengang wirklich errathen haben könne, wollen wir gern zugeben; entscheiden wird sich aber hierüber um so weniger etwas lassen, da dergleichen Untersuchungen auf sehr mannigfaltige Art angegriffen werden können, und überdies auch sonst Spuren vorhanden sind, dass der grosse Grieche im Besitz mancher nichts weniger als gemeiner Wahrheiten und Kunstgriffe, selbst aus der höhern Arithmetik, gewesen sein muss.

Eine Frage bleibt übrigens hier noch übrig, warum nemlich \textsc{Archimed}, wenn er seine genäherten Werthe methodisch gefunden hat. bei den grössern bis zum eilften Gliede gegangen ist, da er doch bei den kleinern nur bis zum achten ging; man sollte glauben, er würde bei jenen sich mit dem neunten Gliede \(\frac{362}{209}\) begnügt haben, welches immer zur Ausmittelung der untern Grenze \(3\frac{10}{71}\) hinreichend gewesen wäre, und könnte vielleicht verleitet werden, hieraus die Folge zu ziehen, dass \textsc{Archimed} doch den Bruch \(\frac{1351}{780}\) durch eine Art von glücklichem Zufall gefunden habe, und der einfachere \(\frac{362}{209}\) ihm entgangen sei. Hr. M. glaubt, \textsc{Archimed} habe jenen Bruch desswegen gewählt, weil er der einfachste von denen sei, deren Zähler zu der Ordnung der Tausender gehören, so wie er den Bruch \(\frac{265}{153}\) als den einfachsten aus der Ordnung der Hunderter gewählt habe: allein dieser Grund scheint uns nicht befriedigend. Wir finden es vielmehr wahrscheinlicher, dass er den Bruch \(\frac{1351}{780}\) desswegen vorzog, weil er fand, dass derselbe zufälliger Weise beim weitern Fortgange der Rechnung eine bequeme Vereinfachung darbietet, so dass sich beim 24 Eck für dasjenige Verhältniss, welches, nach unsrer Art zu reden, \(1:\operatorname{cotang}7^0\,30'\) ist, eine äusserst nahe Grenze sehr
